% !TEX root = ../../ctfp-print.tex

\lettrine[lhang=0.17]{我}{意识到我们可能正逐渐远离编程}，深入到艰深的数学领域。但你永远无法预知下一次编程领域的重大革命会带来什么，也无法知道理解它所需的数学知识是什么。目前有一些非常有趣的想法正在传播，例如带有连续时间特性的函数式响应式编程、通过依赖类型扩展Haskell类型系统，或者在编程中探索同伦类型论。

迄今为止，我一直随意地将类型等同于值的\emph{集合}。这种做法并不严格正确，因为这种方法没有考虑到：在编程中我们是通过\emph{计算}来获得值的，而计算是一个耗时的过程，在极端情况下甚至可能不会终止。发散性计算（Divergent computations）是每个图灵完备语言的固有特性。

从基础上看，集合论可能不适合作为计算机科学甚至数学本身的基石也有其原因。一个恰当的类比是：集合论就像绑定特定架构的\emph{汇编语言}。如果你想让你的数学理论运行在不同的架构上，就需要使用更通用的工具。

一种可能性是用\emph{空间}替代集合。空间具有更多结构，并且可以不依赖集合来定义。通常与空间相关联的是拓扑学（topology），这是定义连续性等概念所必需的。而传统的拓扑学方法正如你所料，是通过集合论实现的。特别是，子集（subset）的概念在拓扑学中处于核心地位。不出所料，范畴论学者已将这一思想推广到除$\Set$以外的其他范畴。具有恰当选取性质、能够替代集合论的这类范畴被称为\newterm{拓扑斯}（topos，复数形式为topoi），它提供了一种广义的子集概念。

\section{子对象分类器}

让我们先尝试用函数而非元素来表达子集的概念。任何从集合$a$到$b$的函数$f$都定义了一个$b$的子集——即$a$在$f$下的像集。但存在许多不同函数定义同一子集的情况。我们需要更精确的定义：首先考虑单射函数（injective functions）——即不会将多个元素压缩为一个元素的函数。单射函数将一个集合"注入"到另一个集合中。对于有限集合，你可以将单射函数可视化为连接两个集合元素的平行箭头。当然，第一个集合不能比第二个集合更大，否则箭头必然汇聚。仍存在一定模糊性：可能存在另一个集合$a'$和另一个单射函数$f'$同样选取相同子集。但你可以轻易验证这样的集合必然是与$a$同构的。我们可以利用这个事实，将子集定义为由定义域同构关系关联的一族单射函数。更精确地说，我们说两个单射函数：
\begin{align*}
  f  & \Colon a \to b  \\
  f' & \Colon a' \to b
\end{align*}
是等价的，当且仅当存在一个同构：
$$h \Colon a \to a'$$
使得：
$$f = f'\ .\ h$$
这样一族等价的单射函数定义了$b$的一个子集。

\begin{figure}[H]
  \centering
  \includegraphics[width=0.4\textwidth]{images/subsetinjection.jpg}
\end{figure}

\noindent
如果我们用单态射（monomorphism）替换单射函数，这个定义可以提升到任意范畴。提醒一下，单态射$m$从$a$到$b$由其泛性质定义：对任意对象$c$和任意一对态射：
\begin{align*}
  g  & \Colon c \to a \\
  g' & \Colon c \to a
\end{align*}
若满足：
$$m\ .\ g = m\ .\ g'$$
则必有$g = g'$。

\begin{figure}[H]
  \centering
  \includegraphics[width=0.4\textwidth]{images/monomorphism.jpg}
\end{figure}

\noindent
在集合范畴中，这个定义更容易理解：若函数$m$不是单态射，则它会将$a$的两个不同元素映射到$b$的单一元素。此时我们可以找到两个仅在这两个元素处不同的函数$g$和$g'$，通过$m$的后复合会掩盖这种差异。

\begin{figure}[H]
  \centering
  \includegraphics[width=0.4\textwidth]{images/notmono.jpg}
\end{figure}

\noindent
定义子集的另一种方式是使用称为特征函数（characteristic function）的单一函数。它是从集合$b$到两元素集合$\Omega$的函数$\chi$。该集合的一个元素被指定为"真"（true），另一个为"假"（false）。这个函数为$b$中属于子集的元素分配"真"，不属于的分配"假"。

还需要明确指定$\Omega$中元素为"真"的含义。我们可以使用标准技巧：用从单元素集合到$\Omega$的函数来表示。我们将这个函数命名为$\mathit{true}$：
$$\mathit{true} \Colon 1 \to \Omega$$

\begin{figure}[H]
  \centering
  \includegraphics[width=0.4\textwidth]{images/true.jpg}
\end{figure}

\noindent
这些定义可以结合起来，不仅定义了子对象，还无需涉及元素即可定义特殊对象$\Omega$。关键在于我们希望态射$\mathit{true}$代表一个"通用"子对象。在$\Set$中，它从两元素集合$\Omega$中选取单元素子集。这已经是最通用的情形了。显然这是一个真子集，因为$\Omega$还有一个元素不在该子集中。

在更一般的场景中，我们定义$\mathit{true}$是从终对象到\emph{分类对象}$\Omega$的单态射。但我们需要定义分类对象本身。我们需要一个泛性质将该对象与特征函数联系起来。事实证明，在$\Set$中，沿特征函数$\chi$对$\mathit{true}$作拉回（pullback）既定义了子集$a$，也定义了将其嵌入$b$的单射函数。以下是拉回图示：

\begin{figure}[H]
  \centering
  \includegraphics[width=0.4\textwidth]{images/pullback.jpg}
\end{figure}

\noindent
分析这个图示。拉回方程为：
$$\mathit{true}\ .\ \mathit{unit} = \chi\ .\ f$$
函数$\mathit{true}\ .\ \mathit{unit}$将$a$的每个元素映射为"真"。因此$f$必须将$a$的所有元素映射到那些使$\chi$为"真"的$b$元素。根据定义，这些正是特征函数$\chi$指定的子集元素。因此$f$的像集确实就是该子集。拉回的泛性质保证了$f$是单射。

这个拉回图示可用于定义除$\Set$外其他范畴的分类对象。这样的范畴必须具有终对象（用于定义单态射$\mathit{true}$），并且必须具有拉回（实际要求是必须具有所有有限极限，拉回是有限极限的一个例子）。在此前提下，我们通过以下性质定义分类对象$\Omega$：对每个单态射$f$，存在唯一的态射$\chi$完成拉回图示。

分析最后这个陈述：构造拉回时我们给定三个对象$\Omega$、$b$和$1$，以及两个态射$\mathit{true}$和$\chi$。存在拉回意味着我们可以找到最佳的对象$a$，配备两个态射$f$和$\mathit{unit}$（后者由终对象定义唯一确定），使得图示可交换。

此处我们解决的是不同的方程组：我们在变化$a$和$b$的同时求解$\Omega$和$\mathit{true}$。对于给定的$a$和$b$可能存在也可能不存在单态射$f \Colon a \to b$。但如果存在，我们希望它是某个$\chi$的拉回。更重要的是，我们希望这个$\chi$由$f$唯一确定。

我们不能说单态射$f$和特征函数$\chi$之间存在一一对应，因为拉回仅在同构意义下唯一。但请回想我们早先将子集定义为等价单射族的方式。可以通过将$b$的子对象定义为等价单态射族来推广这个概念。这个单态射族与我们的图示等价拉回族之间存在一一对应。

因此我们可以将$b$的子对象集$\mathit{Sub}(b)$定义为单态射族，并发现它与从$b$到$\Omega$的态射集同构：
$$\mathit{Sub}(b) \cong \cat{C}(b, \Omega)$$
这恰好是两个函子的自然同构。换句话说，$\mathit{Sub}(-)$是一个可表（逆变）函子，其表示对象为$\Omega$。

\section{拓扑斯}

一个拓扑斯是满足以下条件的范畴：

\begin{enumerate}
  \tightlist
  \item
        是笛卡尔闭范畴：具有所有积、终对象以及指数对象（定义为积的右伴随函子），
  \item
        存在所有有限图式的极限，
  \item
        具有子对象分类器 $\Omega$。
\end{enumerate}

这组性质使得拓扑斯在大多数应用中成为集合范畴 $\Set$ 的自然替代。它还具有由定义推导出的额外性质。例如，拓扑斯具有所有有限余极限，包括始对象。

将子对象分类器定义为两个终对象的余积（和）似乎很直观——这正是它在 $\Set$ 中的表现形式——但我们希望更一般的框架。满足这一条件的拓扑斯被称为布尔拓扑斯。

\section{拓扑斯与逻辑}

在集合论中，特征函数可被解释为定义集合元素性质的一种方式——即对某些元素成立而对另一些元素不成立的\newterm{谓词(predicate)}。谓词$\mathit{isEven}$从自然数集中选出偶数子集。在拓扑斯中，我们可以将谓词概念推广为从对象$a$到$\Omega$的态射。这就是为什么$\Omega$有时被称为真值对象。

谓词是逻辑的基本构建模块。拓扑斯包含研究逻辑所需的所有机制：其乘积对应逻辑合取（逻辑\emph{且}），余积对应析取（逻辑\emph{或}），指数对象对应蕴含式。除排中律（或等价的双重否定消除律）外，所有标准逻辑公理都在拓扑斯中成立。因此拓扑斯的逻辑对应于构造性逻辑或直觉主义逻辑。

直觉主义逻辑一直在稳步发展，并意外地获得了计算机科学的支持。经典的排中律基于绝对真理的信念：任何命题要么为真要么为假，正如古罗马人所说，\emph{tertium non datur}（不存在第三种可能）。但我们判断某命题真假的唯一方式是能否证明或证伪它。证明是一个过程、一种计算——而我们知道计算需要时间和资源，在某些情况下甚至可能永不终止。若某个命题无法在有限时间内被证明，就断言其为真并无实际意义。具有更精细真值对象的拓扑斯为建模有趣的逻辑提供了更通用的框架。

\section{挑战}

\begin{enumerate}
  \tightlist
  \item
        证明函数 $f$ 是特征函数沿 $\mathit{true}$ 的拉回，该函数必然是单射。
\end{enumerate}